\chapter*{摘要}
\addcontentsline{toc}{chapter}{摘要}
大模型智能体在读取网页、工具返回和长期记忆时，可能把其中的恶意指令当作任务要求，进而读取秘密、外发数据或推送代码。安全研判又需要保留并分析危险原文，单纯删除或拒绝可疑内容会损失业务结果。

溯界在智能体宿主与工具之间核验实际动作，以可信任务契约确定资源和副作用范围，以来源用途信封保留不可信数据的传播限制，以放行、审批、阻断三态控制工具执行。系统支持真实MCP调用、跨角色委派、跨会话记忆、动作绑定审批，以及来源图和实际影响审计。

作品已实现13个开发业务入口。M25同事实脚本比较包含48个单元，46个完成合法业务，两个对照单元发生禁止的Git推送；14个来源代表单元均完成业务，其中三个不安全配置发生禁止写入，对应完整配置阻断。集成版本本地回归记录为1025项通过。上述为开发期确定性校准；同版真实模型主比较与未见任务评测仍待完成。

\noindent\textbf{关键词：}大模型智能体；提示注入；工具授权；MCP；来源用途约束
\clearpage
